\scanpage{31}

La condition a) permet déjà de définir le
$\pi_1$ relatif
\[
\pi_1(\sigma)
 = \pi_1(\Pi_{D^*}/\Pi_D),
\]
qui est un syst. local de gr. \unsure{pro|pr.} abéliens sur $\Pi_D$
\inserted{\unsure{i.e.} un foncteur
$(\Pi_D)^{\mathrm{o}}\to\mathrm{Ens}$,}
et la dernière donnée : \deleted{est un}

$4^\circ)$ Un isom. kummérien
\[
\deleted{\illegible{1}}\,k:\quad
\pi_1(\sigma)\simeq
\psi^*(\mathcal T)
\]

Si on a deux systèmes de cette nature
\[
\Pi=(\Pi_U,\Pi_D,\Pi_{D^*},\varphi,\psi,\rho,\sigma,\alpha,k)
\]
et $\Pi'=(\Pi'_U,\Pi'_D,\ldots)$, un \uline{morphisme}
de $\Pi'$ dans $\Pi$ est un système de $3$ foncteurs
\[
\begin{gathered}
f_U:\Pi'_U\to\Pi_U\\
f_D:\Pi'_D\to\Pi_D\\
\deleted{\illegible{1}}f_{D^*}:\Pi'_{D^*}\to\Pi_{D^*}
\end{gathered}
\]
et de \inserted{quatre}\deleted{\illegible{1}}
isomorphismes de commutation
$\unsure{\alpha_{D^*D^*}|\alpha_{D^*D}},\alpha_{D^*U},\alpha_{Ue},\alpha_{De}$
pour les quatre faces du prisme
\diagram{
\begin{tikzcd}[row sep=large,column sep=large]
 & \Pi'_{D^*}\arrow[rr,"\sigma'"]\arrow[dd,"\unsure{\rho'|\varphi'}"]
   \arrow[dl,"f_{D^*}"'] &&
   \Pi'_D\arrow[dd,"\psi'"]\arrow[dl,"f_D"']\\
 \Pi_{D^*}\arrow[rr,"\sigma",near start]\arrow[dd,"\rho"'] &&
 \Pi_D\arrow[dd,"\psi",near end]\\
 & \Pi'_U\arrow[rr,"\varphi'"]\arrow[dl,"f_U"'] &&
 \Pi_e\arrow[dl,"="']\\
 \Pi_U\arrow[rr,"\varphi"'] && \Pi_e
\end{tikzcd}
}
\deleted{\illegible{2}}

\scanpage{32}
\mspage{14}

satisfaisant une \inserted{équation} de compatibilité avec $\alpha,\alpha'$ que je n'écris pas
--- signifiant que les \uline{deux}
\deleted{\illegible{1}} isomorphismes \inserted{$u,v$} de foncteurs
\diagram{[Deux flèches parallèles de $\Pi'_{D^*}$ vers $\Pi_e$.
La flèche supérieure porte $\varphi\circ\rho\circ f_{D^*}$ ;
la flèche inférieure porte $\unsure{\mathrm{id}}\circ\psi'\circ\sigma'$.
Entre les deux, deux flèches descendantes, notées $u$ et $v$.]}
obtenus respectivement, $u$ en utilisant successivement
$\alpha,\alpha_{D^*D},\alpha_{De}$, \deleted{\unsure{et}} $v$ en
utilisant successivement $\alpha_{D^*U},\alpha_{Ue},\alpha'$, sont égaux.
[Cette compatibilité pourrait s'exprimer en interprétant la donnée
de $\Pi$ comme celle d'une catégorie
fibrée \inserted{$\Pi$} sur la \deleted{\unsure{de diagramme en forme}}
\inserted{``catégorie} carré''
\[
Q:
\]
\diagram{[Losange de sommets $D^*$ en haut, $\unsure{U}$ à gauche,
$D$ à droite, $e$ en bas. Le côté $D^*\to U$ porte $\rho_0$ ;
les trois autres côtés portent $\sigma_0$ entre $D^*$ et $D$,
$\varphi_0$ entre $U$ et $e$, $\psi_0$ entre $D$ et $e$.
Les pointes de ces trois traits ne sont pas nettement discernables.]}
où $\varphi_0\rho_0=\psi_0\sigma_0$.

à restriction à $\{e\}$ imposée, et en prenant des foncteurs
\inserted{cartésiens} entre catégories fibrées\ldots].
De plus, on exige une autre compatibilité, savoir que

\scanpage{33}

\deleted{l'hom.} l'hom. de syst. locaux en groupes abéliens sur $\Pi'_D$
\[
(f_D)^*\bigl(\pi_1(\sigma)\bigr)
\longleftarrow \pi_1(\sigma')
\]
défini à l'aide de : $f_{D^*},f_D,\alpha_{D^*D}$
rende commutatif le diagramme suivant
d'\deleted{\illegible{1}}isomorphismes \inserted{de}
\deleted{foncteurs} systèmes locaux sur $\Pi'_D$
\diagram{
\begin{tikzcd}[row sep=large,column sep=large]
(f_{D^*})^*\bigl(\pi_1(\sigma)\bigr)
 & \pi_1(\sigma')\arrow[l]\\
(f_{D^*})^*\bigl(\psi^*(\mathcal T)\bigr)
 \arrow[u,"(f_{D^*})^*(k)","\sim"']
 \arrow[d,"\mathrm{can}","\sim"']
 & \psi'^*(\mathcal T)
 \arrow[u,"k'","\sim"']\arrow[dl,"\sim"']\\
(\psi f_D)^*(\mathcal T) &
\end{tikzcd}
}
déduit de $\alpha_{De}:\psi'\simeq\psi\circ f_D$
\deleted{\illegible{1}}

\margin{\uline{non}, \unsure{ça|ce} ne marche que pour les
cas de morphismes étales, sinon il faut faire
intervenir la multiplication par des
``degrés de ramification'' $d_{i'}$ $(i'\in\pi_0(\Pi'_{D^*}))$.}

Pour $\Pi,\unsure{\Pi^*|\Pi'}$ fixes, les systèmes
\[
(f,\alpha):=(f_U,f_D,f_{D^*},\alpha_{D^*D},\alpha_{D^*U},
\alpha_{U,e},\alpha_{D,e})
\]
précédents forment une catégorie de \unsure{faç.} naturelle
--- en fait un groupoïde --- en prenant comme morphismes
\inserted{$\mu$} de $(f',\alpha')$ dans $(f,\alpha)$
les triplets des morphismes
(foncteurs profinis)
\[
f'_U\xrightarrow{\mu_U}\unsure{f_U},\qquad
f'_D\xrightarrow{\mu_D}f_D,\qquad
f'_{D^*}\xrightarrow{\mu_{D^*}}f_D
\]

\scanpage{34}
\mspage{15}

satisfaisant \deleted{des} quatre conditions de compatibilité avec
$\alpha_{D^*D}$ et $\alpha'_{D^*D}$, avec $\alpha_{D^*U}$ et
$\alpha'_{D^*U}$, avec $\alpha_{Ue}$ et $\alpha'_{U,e}$, avec
$\alpha_{D,e}$ et $\alpha'_{D,e}$ respectivement.
(i.e. on travaille sur une ss-catégorie pleine de la
\unsure{$2$} cat. des \uline{$\mathrm{Hom}$} entre catégories
fibrées sur $Q$, à fibre en $e$ fixée\ldots).
J'ai l'impression que ce groupoïde
\uline{$\mathrm{Hom}$}$((f',\alpha'),(f,\alpha))$ est
tjrs \uline{rigide}, i.e.
$\mathrm{Aut}(f,\alpha)$ est toujours réduit au groupe unité
--- j'ai la flemme de vérifier --- donc que si $(f',\alpha')$ et
$(f,\alpha)$ sont isomorphes, l'isom. en question est unique.
Quoi qu'il en soit, on posera

$\mathrm{Hom}((f',\alpha'),(f,\alpha))=\pi_0$
\uline{$\mathrm{Hom}$}$((f',\alpha'),(f,\alpha))$

D'où une \uline{catégorie} des
\deleted{couples \unsure{$(X,S)$}} systèmes
$(\Pi_U,\Pi_D,\Pi_{D^*},\varphi,\psi,\rho,\sigma,\alpha,
\deleted{\illegible{1}}k)$

\scanpage{35}

On a un foncteur des couples $(X,S)$
(\unsure{i.e.|où} $X$ schéma \inserted{loc.} de t.f. sur $K$,
$S$ \deleted{partie}\inserted{sous-schéma}
\inserted{fermé \deleted{\illegible{1}}} \deleted{\unsure{discrète|discret}} de $X$
\inserted{étale sur $K$} (\unsure{telle|tel} que $\forall s\in S$, $X$
soit lisse de dim relative $1$ en $s$)
vers cette catégorie bordélique $\mathcal B$.

\margin{\unsure{Morphismes|Les morphismes} $(X',S')\to(X,S)$ sont les
morphismes $f:X'\to X$ tels que $f^{-1}(S)_{\text{réd}}=S'$.}

\uline{Conjecture bordélique 7.1.}
Quand on se borne aux $(X,S)$ tels que $X$ projectif lisse de dim
$\leq 1$, et qui de plus sont anabéliens, alors le foncteur
précédent est pleinement fidèle.

\margin{\uline{NB} La fidélité est facile\ldots}

Description de la catégorie bordélique en termes de théorie des groupes :

Soit $I=\pi_0(\Pi_{D^*})\simeq\pi_0(\Pi_D)$.
Choisissons un \unsure{élément|élt} \inserted{$D_i^*$} dans chaque composante de
$\Pi_{D^*}$, et soit $D_i=\sigma(D_i^*)$.
Pour tt $i$, choisissons un isomorphisme
\[
\boxed{\psi(D_i)\xrightarrow[\simeq]{\unsure{\partial_{D_i}|\lambda_{D_i}}}
\mathrm{Spec}(\overline K)}.
\]
\deleted{\unsure{d'où} \illegible{3}}
Quitte à remplacer l'objet par un ``sous-objet'' isomorphe,

\scanpage{36}
\mspage{16}

OPS que $\Pi_{D^*}$ est la catégorie somme des catégories
$[\mathcal E_i]$ définies par les
$\mathcal E_i=\mathrm{Aut}(D_i^*)$, $\Pi_D$
\deleted{est} \deleted{\illegible{1}}\inserted{la catégorie}
somme de catégories $[\Gamma_i]$ définies par les
$\Gamma_i=\mathrm{Aut}(D_i)$, le foncteur $\sigma$ s'exprimant
par un système d'homs
\[
\sigma_i:\mathcal E_i\to\Gamma_i\qquad(i\in I)
\]
qui sont surjectifs de noyaux abéliens,
\inserted{le foncteur $\psi$ par un système d'\uline{inclusions}
$\psi_i:\Gamma_i\subset\Gamma$}
et la donnée de $k$ équivaut en fait à des isom.
\[
k_i:\mathrm{Ker}\,\sigma_i\simeq T(\overline K)
\]
compatibles avec les op. de $\Gamma_i$ et $\Gamma$ sur les deux
\unsure{mbs|mbres} respectivement et les inclusions $\psi_i$.

\margin{$^*$Re N.B. Deux bases différentes \unsure{définissent}
des scindages différents!}

On peut dire que la donnée de $(\sigma,\psi,k)$ est exprimée
la donnée du syst. $(\mathcal E_i)_{i\in I}$, d'un système
\unsure{de} \deleted{homs} suites exactes
\deleted{$p_i:\mathcal E_i\to\Gamma$}
\deleted{\illegible{5} et des isom. \illegible{1}}
\[
\boxed{1\to T(\overline K)\xrightarrow{k_i}\mathcal E_i
\xrightarrow{p_i}\Gamma}
\]
telles que les $\Sigma_i
=\mathrm{Im}\,p_i$ soient ouverts, et que $k_i$ soit compatible
avec les opérations de $\Sigma_i\simeq\mathrm{Coker}\,k_i$
(i.e. de $\mathcal E_i$ et $\Gamma$ sur les deux termes respectivement).

\margin{\uline{NB} Les extensions des $\Gamma_i$ par
$T(\overline K)$ obtenues par des situations géométriques
splittent --- \illegible{1} le choix d'une uniformisante en $s_i$
définit un splittage, \unsure{et même} \illegible{1} le choix
d'une base de l'espace tangent en $s$\ldots}

\scanpage{37}

Supposant d'autre part (pour simplifier)
\deleted{\unsure{$X$}}\inserted{$\Pi_U$} connexe,
\deleted{\unsure{i.e.} $U=X\setminus S$ connexe, et}
et choisissant un élt $\widetilde U$ de $\Pi_U$ et un isom
\deleted{entre $\varphi(\widetilde U)$ et $\mathrm{Spec}\,K$}
\[
\varphi(\widetilde U)\xrightarrow{\unsure{\partial_U|\lambda_U}}
\mathrm{Spec}(\overline K)
\]
donc (quitte à remplacer $\Pi_U$ par un groupoïde équivalent)
OPS $\Pi_U$ réduit à $\widetilde U$, et $\Pi_U$ est donné alors par
un groupe $\mathcal E$, et $\varphi$ par un hom de groupes
\[
\boxed{\mathcal E
\xrightarrow{\deleted{\illegible{1}}\,\unsure{\varphi_{\widetilde U,\partial_U}|\varphi_{\widetilde U,\lambda_U}}}
\Gamma}
\]
\inserted{(ou $p$)}
dont l'image $\Sigma\subset\Gamma$ est donc un ss-groupe ouvert
de $\Gamma$, et le noyau \unsure{sera noté} $\unsure{\pi|\Pi}$
\[
\boxed{1\to\unsure{\pi|\Pi}\xrightarrow{\unsure{\kappa|\iota}}\mathcal E
\xrightarrow{\unsure{\varphi|p}}\Sigma_{\unsure{0}}\to1}
\]
\diagram{[Sous $\Sigma_{\unsure{0}}$, un signe d'inclusion vertical vers $\Gamma$.]}

\margin{\unsure{Bien noté} Dans les situations géom.
\uline{NB} (Cette extension de noyaux de groupes splitte, i.e.
$\exists$ ss-groupe \unsure{ouvert} $\Sigma_0$ \unsure{dans} $\Gamma$,
\deleted{\unsure{d'indice} \illegible{1}} \inserted{et un} relèvement
$\Sigma_0\to\mathcal E$\ldots}

[Ayant remplacé le $\Pi_U$ initial par une ss-catégorie pleine
plus petite, on sera obligé de modifier $\rho$ à isom près,
pratiquement en choisissant pour chaque $i\in I$ un isom.
\[
\rho(D_i)\xrightarrow{\mu_i}\widetilde U].
\]
\deleted{\illegible{1}}

\scanpage{38}
\mspage{17}

moyennant quoi $\rho$ s'explicite simplement par des
homomorphismes de groupes
\[
\boxed{\rho_i:\mathcal E_i\longrightarrow\mathcal E}
\]
Il reste à expliciter l'isom de commutation
\[
\alpha:\varphi\rho\to\psi\sigma
\]
qui est défini par un système d'éléments
\[
\boxed{\gamma_i\in\Gamma\qquad(i\in I)}
\]
tels que
\[
\unsure{(\alpha)|(*)}\qquad
\varphi\circ\rho_i=\mathrm{int}(\gamma_i)\circ p_i.
\]
ce qui implique d'ailleurs que $\rho_i$ \deleted{\illegible{1}}
applique $\mathrm{Ker}\,p_i$ dans $\mathrm{Ker}\,p$, i.e. induit
un hom de suites exactes
\diagram{
\begin{tikzcd}[row sep=large,column sep=large]
1\arrow[r] & T(\overline K)\arrow[r]\arrow[d] &
\mathcal E_i\arrow[r]\arrow[d,"\rho_i"] &
\deleted{\unsure{\mathcal E_i}}\inserted{\Sigma_i}
\arrow[r]\arrow[d,hook] & 1\\
1\arrow[r] & \unsure{\pi|\Pi}\arrow[r] & \mathcal E\arrow[r] &
\deleted{\unsure{\mathcal E_0}}\inserted{\Sigma_0}
\arrow[r] & 1
\end{tikzcd}
[À droite de $\Sigma_i$, en exposant, $\subset\Gamma$ ; sous $\Sigma_0$, un signe d'inclusion et $\Gamma$.]
}
avec un hom induit $\Gamma_i\to\Gamma_0$ injectif
--- et ceci posé, la relation de compatibilité
\unsure{$(\alpha)$|$(*)$} devient une relation sur des
monomorphismes de groupes
\diagram{
\begin{tikzcd}
\Gamma_i\arrow[rr,hook]\arrow[dr] & & \Gamma\\
 & \Gamma_0\arrow[r,hook] & \Gamma
\end{tikzcd}
}

\scanpage{39}

\deleted{\uline{NB} Quand $I$ est vide i.e. $\Pi_{D^*},\Pi_D$
vides, alors les systèmes ``bordéliques'' \unsure{découlant}
de situations géom. correspondent à : $\unsure{\pi|\Pi}=\{1\}$, et
$\mathcal E$ se réduit à un \unsure{ss-groupe}}

Ainsi, on a une description relativement simple des systèmes
bordéliques ``réduits'' (i.e. où dans les groupoïdes
$\Pi_U,\Pi_{D^*},\Pi_D$, chaque composante connexe a exactement
un objet, et où \inserted{de plus} on force en quelque sorte
$\unsure{\Pi_1(K)|\Pi_e(K)}$ à n'avoir que l'objet
$\mathrm{Spec}(\overline K)$.
Mais on est retrouvé \unsure{en tournant}
\unsure{pr|par} la détermination des \uline{morphismes} des systèmes
\[
\mathcal G=(\inserted{I,}G_i,p_i,k_i,\mathcal E,\varphi,\rho_i,
\gamma_i)
\quad\hbox{et}\quad
\mathcal G'=(I',G'_{i'},\ldots),
\]
\inserted{(ou $p$?)}
\deleted{\illegible{2}} disons dans le \inserted{sens}
\[
f:\mathcal G'\to\mathcal G.
\]
Il faut donc \inserted{(pour $f_{D^*}$)} une application
\[
\boxed{\deleted{\illegible{1}}\,\unsure{\tau=\tau_f}:I'\to I},
\]
et $\forall\unsure{i'|i}$ un hom
\[
\boxed{G'_{i'}\xrightarrow{f_{i'}}G_{\tau i'}},
\]
\deleted{induisant \unsure{par} \illegible{1}}
\inserted{(compte tenu de $f_D$)}
induisant \inserted{\unsure{par passage au}} quotient des homs
\[
\Gamma'_{i'}\xrightarrow{g_{i'}}\Gamma_{\tau i'}
\]
\deleted{\unsure{mais}}

\margin{Les ``objets bordéliques simplifiés'' sont donc les
systèmes d'hom. de suites exactes
\diagram{
\begin{tikzcd}[row sep=large,column sep=large]
1\arrow[r] & T(\overline K)\arrow[r,"k_i"]
\arrow[d,"\unsure{\rho_i^0}"] &
\mathcal E_i\arrow[r,"p_i"]\arrow[d,"\rho_i"] &
\Gamma\arrow[d,"\mathrm{int}(\gamma_i)"]\\
1\arrow[r] & \unsure{\pi|\Pi}\arrow[r,"\unsure{\kappa}"] &
\mathcal E\arrow[r,"p"] & \Gamma
\end{tikzcd}
}}

\scanpage{40}
\mspage{18}

de ss-groupes ouverts de $\Gamma$,
et la donnée \inserted{$f_{D^*},f_D,$}
$\alpha_{D^*,D},\alpha_{D,e}$ équivaut alors
à la donnée d'éléments $\alpha_i$
$(\unsure{i\in I|i'\in I'})$ de $\Gamma$,
\inserted{$\boxed{\alpha_{i'}\in\Gamma}
\quad(i'\in I')$}
tels que
\deleted{$\forall i\in I$, $g\in\Gamma_i$, \illegible{2}}
\[
g_{i'}(\gamma')=\mathrm{int}(\alpha_{i'})(\gamma')
\qquad\forall i'\in I',\quad\gamma'\in\Gamma'_{i'}.
\]
La donnée de $f_U$ équivaut à la donnée d'un hom de groupes
\[
\boxed{f_{\mathcal E}:\mathcal E'\to\mathcal E}
\]
celle de $\alpha_{D,e}$ équivaut à la donnée de
\[
\boxed{\alpha\in\Gamma}
\]
tel que
\[
\unsure{(\alpha)|(*)}\qquad
\varphi f_{\mathcal E}=\mathrm{int}(\alpha)\varphi
\]
de sorte que $f_{\mathcal E}$ induit un hom de suites exactes
\diagram{
\begin{tikzcd}[row sep=large,column sep=large]
1\arrow[r] & \unsure{\pi|\Pi}'\arrow[r]\arrow[d] &
\mathcal E'\arrow[r]\arrow[d,"f_{\mathcal E}"] &
\Gamma'_0\arrow[r]\arrow[d,"\unsure{f_{\Gamma_0}|f_{\Sigma_0}}"] & 1\\
1\arrow[r] & \unsure{\pi|\Pi}\arrow[r] & \mathcal E\arrow[r] &
\Gamma_0\arrow[r] & 1
\end{tikzcd}
[Au-dessus de $\Gamma'_0$, $\Gamma$ avec un signe d'inclusion
vertical ; sous $\Gamma_0$, un signe d'inclusion et $\Gamma$.]
}
et moyennant cela, la condition dite \unsure{$(\alpha)$|$(*)$}
\deleted{est}\inserted{devient} une condition sur des
\inserted{inclusions de} ss-groupes de $\Gamma$ :
\[
f_{\Gamma_0}(\gamma')=\mathrm{int}(\alpha)\gamma'
\quad\hbox{si}\quad\gamma'\in\Gamma'_0.
\]
Il faut expliciter encore (en plus des
$f_{D^*},f_D,f_U,\alpha_{D^*D},\alpha_{De},\alpha_{Ue}$
déjà explicités) le
